\documentclass[12pt]{article}
\usepackage[utf8]{inputenc}
\usepackage[T1]{fontenc}
\usepackage{hyperref}
\title{Tupinikin Language Resources for K-12 Curriculum Development}
\author{Compiled by Claude Sonnet 5 (claude-sonnet-5) for Andr\'as Kornai}
\date{}
\begin{document}
\maketitle
\section{Introduction}
Tupinikin (ISO 639-3 tpk) is a dormant Tupian language of Brazil, associated with the Tupinikim people of Esp\'irito Santo. It has no first-language speakers. Searches found academic theses on revitalisation, a university teacher-training course and a state cultural project. They found no published textbooks or dictionaries. The people's name is also spelled Tupiniquim.
\section{Primary Educational Resources}
No published textbooks were found. The UFES Intercultural Indigenous teaching degree (Tupinikim and Guarani) includes a course on producing dictionaries and teaching materials in Tupi-Guarani languages: \url{https://indigena.ufes.br/sites/indigena.ufes.br/files/field/anexo/plano_de_ensino_ensino_e_pesquisa_em_linguas_indigenas.docx.pdf}
\section{Linguistic and Reference Materials}
\begin{itemize}
\item UFES dissertation on Tupi in the Tupinikim village of Caieiras Velha, Aracruz, as a question of language policy: \url{https://repositorio.ufes.br/items/5c3d17a4-0ed7-4d05-98d7-649448073a39/full}
\item UFRJ dissertation on a language policy built by the Tupinikim people: \url{https://pantheon.ufrj.br/handle/11422/24447}
\item UnB thesis (Valdemir de Almeida Silva, 2016): \url{https://repositorio.unb.br/bitstream/10482/22777/1/2016_ValdemirdeAlmeidaSilva.pdf}
\end{itemize}
\section{Religious and Community Resources}
None found in search results.
\section{Digital and Audio Resources}
\begin{itemize}
\item ``L\'ingua Viva'' project of the Esp\'irito Santo government. It has a digital collection of indigenous artists' work and a podcast, with emphasis on the mother tongue: \url{https://www.es.gov.br/Noticia/projeto-destaca-artes-indigenas-como-forma-de-preservacao-da-lingua-tupiniquim}
\end{itemize}
\section{Curriculum Development Context}
The Caieiras Velha dissertation reports that the community is building a revitalisation project. It finds that legal guarantees of indigenous-language education exist but are hard to implement, and that school teaching alone will not secure revitalisation. Another UFES document looks at Tupinikim handicrafts as teaching material in early childhood education: \url{https://indigena.ufes.br/sites/indigena.ufes.br/files/field/anexo/tanga_colar_lanca_cocar_e_sambura.pdf}
\end{document}